% =============================================================================
% EUCASS 2027 Paper Draft (Library / Tool Paper)
% "Aetherion: An Open-Source C++23 Library for Structure-Preserving
%  Six-Degree-of-Freedom Flight Dynamics"
%
% Author: Onur Tuncer, Istanbul Technical University
%
% This is the broad "introduce the library" paper.  The companion paper
% (papers/eucass/main.tex) treats the SE(3) RKMK integrator in depth.
%
% NOTE: formatted to match the EUCASS proceedings style.  When the official
% EUCASS 2027 LaTeX class is released, replace
%   \documentclass[a4paper,10pt]{article}
% with
%   \documentclass[twoside]{eucass}
% and remove the manual geometry/font settings below.
% =============================================================================

\documentclass[a4paper,10pt]{article}

% ---------------------------------------------------------------------------
% Packages
% ---------------------------------------------------------------------------
\usepackage[a4paper, top=25mm, bottom=25mm, left=20mm, right=20mm]{geometry}
\usepackage{times}
\usepackage{amsmath, amssymb, amsthm}
\usepackage{mathtools}
\usepackage{bm}
\usepackage{graphicx}
\usepackage{booktabs}
\usepackage{multirow}
\usepackage{hyperref}
\usepackage{xcolor}
\usepackage{caption}
\usepackage{subcaption}
\usepackage{microtype}
\usepackage{natbib}
\usepackage{fancyhdr}
\usepackage{titlesec}
\usepackage{enumitem}
\usepackage{siunitx}
\usepackage{listings}

% ---------------------------------------------------------------------------
% EUCASS-style formatting
% ---------------------------------------------------------------------------
\pagestyle{fancy}
\fancyhf{}
\renewcommand{\headrulewidth}{0.4pt}
\fancyhead[L]{\small\itshape 13th European Conference for Aeronautics and Space Sciences (EUCASS)}
\fancyhead[R]{\small\thepage}
\fancyfoot{}

\titleformat{\section}{\normalsize\bfseries\uppercase}{\thesection.}{0.5em}{}
\titleformat{\subsection}{\normalsize\bfseries}{\thesubsection}{0.5em}{}
\titleformat{\subsubsection}{\normalsize\itshape}{\thesubsubsection}{0.5em}{}

\setlength{\parindent}{0pt}
\setlength{\parskip}{6pt}

% ---------------------------------------------------------------------------
% Listings style for C++
% ---------------------------------------------------------------------------
\definecolor{lstkw}{HTML}{0F4C81}
\definecolor{lstcm}{HTML}{6B7280}
\definecolor{lstst}{HTML}{9A3412}
\definecolor{lstbg}{HTML}{F7F7F9}

\lstdefinestyle{cpp}{
  language=C++,
  basicstyle=\ttfamily\scriptsize,
  keywordstyle=\color{lstkw}\bfseries,
  commentstyle=\color{lstcm}\itshape,
  stringstyle=\color{lstst},
  backgroundcolor=\color{lstbg},
  frame=single,
  rulecolor=\color{lstcm!40},
  framesep=4pt,
  showstringspaces=false,
  columns=fullflexible,
  keepspaces=true,
  breaklines=true,
  captionpos=b,
  morekeywords={concept,requires,constexpr,consteval,noexcept,nullptr,auto,using,template,typename,struct,class,namespace,co_await}
}
\lstset{style=cpp}

% ---------------------------------------------------------------------------
% Math macros
% ---------------------------------------------------------------------------
\newcommand{\SO}{\mathrm{SO}(3)}
\newcommand{\SE}{\mathrm{SE}(3)}
\newcommand{\so}{\mathfrak{so}(3)}
\newcommand{\se}{\mathfrak{se}(3)}
\newcommand{\Exp}{\mathrm{Exp}}
\newcommand{\dexp}{\mathrm{d}\exp}
\newcommand{\ad}{\mathrm{ad}}
\newcommand{\R}{\mathbb{R}}
\newcommand{\norm}[1]{\left\lVert #1 \right\rVert}
\newcommand{\nuB}{\boldsymbol{\nu}_B}
\newcommand{\IB}{\mathbf{I}_B}
\newcommand{\wext}{\mathbf{w}_{\mathrm{ext}}}
\newcommand{\code}[1]{\texttt{\small #1}}

\theoremstyle{definition}
\newtheorem{remark}{Remark}

% ---------------------------------------------------------------------------
% Document
% ---------------------------------------------------------------------------
\begin{document}

% ---------------------------------------------------------------------------
% Title block
% ---------------------------------------------------------------------------
\begin{center}
  {\Large\bfseries
    Aetherion: An Open-Source C++23 Library for\\[4pt]
    Structure-Preserving Six-Degree-of-Freedom Flight Dynamics}

  \vspace{10pt}

  {\normalsize\bfseries Onur Tuncer$^{1}$}

  \vspace{4pt}

  {\small
    $^{1}$Department of Aerospace Engineering, Istanbul Technical University,\\
    Maslak, 34469 Istanbul, Turkey\\
    \texttt{otuncer@itu.edu.tr}}

  \vspace{10pt}
\end{center}

% ---------------------------------------------------------------------------
% Abstract
% ---------------------------------------------------------------------------
\begin{center}
  \begin{minipage}{0.92\linewidth}
    \small\textbf{Abstract}

    Six-degree-of-freedom (6-DOF) flight simulation underpins vehicle design,
    control law verification, and trajectory optimisation, yet the software
    used in practice is typically either proprietary, tied to a single vehicle
    class, or written in a style that resists differentiation and reuse.  This
    paper introduces \emph{Aetherion}, an open-source (MIT) C++23 library for
    atmospheric and space flight vehicle dynamics built on three design
    commitments: Featherstone spatial vector algebra as the sole kinematic and
    dynamic vocabulary; a state that lives explicitly on the product manifold
    $\SE \times \R^{7}$ and is advanced by Runge--Kutta--Munthe-Kaas (RKMK)
    integrators that never leave the group; and scalar-type genericity
    throughout, so that every model evaluates identically under
    \code{double} and under an algorithmic-differentiation (AD) type.  Vehicle
    models are assembled from gravity, aerodynamic, propulsion, and mass
    policies that are checked at compile time by C++20 concepts instantiated
    for \emph{both} scalar types, so a model that is not AD-safe fails to
    compile rather than failing inside a differentiation tape.  The library is
    verified against sixteen check cases from NASA
    TM-2015-218675, spanning ballistic point masses, tumbling rigid bodies,
    wind-shear fields, closed-loop F-16 manoeuvres, polar and equatorial
    circumnavigation, and a two-stage rocket ascent to
    \SI{232}{\kilo\metre}; agreement with the NASA reference is sub-metre for
    the analytic cases and below \SI{1}{\percent} for the full launch-vehicle
    trajectory.  We report the verified convergence order of the integrators,
    describe the extension path for new vehicles, and summarise the continuous
    integration gates---static analysis, sanitizers, include hygiene, and
    complexity metrics---that the code base is held to.  Plants export to
    FMI 2.0/3.0 functional mock-up units, allowing Aetherion models to be
    co-simulated inside third-party environments.
  \end{minipage}
\end{center}

\vspace{6pt}

% ---------------------------------------------------------------------------
\section{Introduction}
\label{sec:intro}
% ---------------------------------------------------------------------------

The equations of motion of a rigid aerospace vehicle are not in dispute.  What
differs between simulation codes---and what determines whether a code can be
reused, differentiated, verified, or trusted at the margins of its envelope---is
how those equations are \emph{expressed} in software.

Three recurring choices in existing 6-DOF codes motivated this work.

\textbf{Flattened state vectors.}  The dominant idiom is to pack attitude,
position, velocity, and mass into a single flat array and hand it to a
general-purpose ODE solver.  This is convenient and wrong in a specific way:
the attitude components are constrained (a unit quaternion, or an orthonormal
matrix) and the solver does not know it.  The constraint is then restored by
periodic re-normalisation, an ad hoc projection applied outside the
integrator that perturbs the trajectory and is invisible to any propagation
Jacobian derived from the solver.

\textbf{Hand-written Jacobians.}  Codes intended for estimation or optimisation
typically carry a second, separately maintained implementation of the dynamics
that returns partial derivatives.  The two implementations drift apart.  The
alternative---finite differences---introduces truncation error precisely where
the Newton or filter update is most sensitive to it.

\textbf{Vehicle-specific plumbing.}  Aerodynamic, propulsion, and mass models
are often hard-wired into the integration loop, so adding a vehicle means
editing the solver rather than supplying a model.

Aetherion is an attempt to remove all three by construction rather than by
convention.  The state is declared on the manifold it actually inhabits and
advanced by an integrator that respects it; every model is written once,
generically over the scalar type, so the differentiated version \emph{is} the
simulated version; and vehicles are assembled from policies whose interfaces
are enforced by the compiler.

\textbf{Contributions.}  This paper contributes:

\begin{enumerate}[leftmargin=*, label=(\arabic*)]
  \item A description of the library architecture, and in particular of the
    compile-time policy/concept mechanism that makes AD-safety a property the
    compiler checks rather than a property the author remembers
    (Section~\ref{sec:design}--\ref{sec:arch}).
  \item A verified numerical core: RKMK integrators on $\SE$ whose empirical
    convergence orders are measured, not assumed, and whose attitude constraint
    residual remains at the round-off level over 600-second trajectories
    (Section~\ref{sec:numerics}).
  \item A verification campaign against sixteen NASA TM-2015-218675 check
    cases with published reference trajectories, reported per scenario
    (Section~\ref{sec:verification}).
  \item An explicit extension path---what a user writes to add a new
    vehicle---together with the interoperability surface (FMI/FMU export,
    DAVE-ML model ingestion, JSON configuration, NASA-compatible CSV output)
    (Sections~\ref{sec:extensibility}--\ref{sec:interop}).
  \item A statement of the engineering process: coding standard, continuous
    integration gates, and test coverage (Section~\ref{sec:quality}).
\end{enumerate}

The library is available under the MIT licence \citep{aetherion_2025}.  The
companion paper \citep{tuncer2027rkmk} treats the implicit SE(3) RKMK
integrator and its Newton/AD machinery in depth; the present paper is
deliberately broader and lighter on derivation.

% ---------------------------------------------------------------------------
\section{Related Software}
\label{sec:related}
% ---------------------------------------------------------------------------

Open-source flight dynamics software is well established but stratified by
purpose.  JSBSim \citep{berndt2004jsbsim} is a mature, widely used
atmospheric flight dynamics model with an extensive XML vehicle description
format, oriented towards flight simulation rather than towards
differentiability or launch-vehicle trajectories.  General multibody
libraries such as RBDL \citep{felis2017rbdl} and Pinocchio
\citep{carpentier2019pinocchio} implement Featherstone's algorithms with high
quality---Pinocchio in particular offers analytical derivatives---but target
robotic articulated systems and do not carry atmosphere, gravity-field, or
launch-vehicle staging models.  On the geometry side, Sophus and manif
\citep{deray2020manif} provide Lie group primitives with Jacobians, but are
libraries of group operations rather than of vehicle dynamics.  MATLAB's
Aerospace Blockset and NASA's Trick/POST families are widely used in
industry but are respectively commercial and export-controlled or
distribution-restricted.

Aetherion occupies the intersection: aerospace vehicle models (atmosphere,
WGS84 gravity, DAVE-ML aerodynamic tables, staging) expressed in Featherstone
spatial algebra, integrated by Lie-group methods, and differentiable
end-to-end.

% ---------------------------------------------------------------------------
\section{Design Principles}
\label{sec:design}
% ---------------------------------------------------------------------------

\subsection{One vocabulary: spatial vector algebra}

All kinematic and dynamic quantities are 6-vectors in Featherstone's spatial
algebra \citep{featherstone2008rbd}: twists $\nuB = [\boldsymbol{\omega}_B^\top,
\mathbf{v}_B^\top]^\top$, wrenches $\mathbf{w} = [\mathbf{m}^\top,
\mathbf{f}^\top]^\top$, the $6\times6$ spatial inertia $\IB$, and the motion
and force cross-product operators $\ad$ and $\ad^*$.  The equation of motion
is then a single line,
\begin{equation}
  \IB(m)\,\dot{\nuB} + \ad^*(\nuB)\,\IB(m)\,\nuB = \wext,
  \label{eq:newton_euler}
\end{equation}
with no separate force and moment bookkeeping and no frame-specific special
cases.  Types are distinct at compile time (\code{Spatial::Twist},
\code{Spatial::Wrench}, \code{Spatial::Inertia}), so a wrench cannot be
silently passed where a twist is expected.

\subsection{The state lives on the manifold}

The 6-DOF state is the triple $(g, \nuB, m)$ on
\begin{equation}
  \mathcal{M} = \SE \times \R^{6} \times \R_{>0},
\end{equation}
represented as a pose $g \in \SE$ plus a Euclidean vector
$x = [\nuB^\top, m]^\top \in \R^{7}$.  The two factors are advanced by
different retractions: $x$ by ordinary weighted sums, $g$ by the Lie
exponential.  Because $g$ is only ever updated as $g \leftarrow g \cdot
\Exp(\Delta)$ with $\Delta \in \se$, the rotation factor is a product of exact
$\SO$ elements and never needs re-orthonormalisation.

\subsection{Scalar-type genericity}

Every model---gravity, atmosphere, aerodynamic table interpolation,
propulsion, the exponential map, the $\dexp^{-1}$ operator, the Newton--Euler
right-hand side---is a template over a scalar type $S$.  Instantiating with
$S = $ \code{double} yields the simulation path; instantiating with
$S = $ \code{CppAD::AD<double>} \citep{cppad} yields a differentiable path
through \emph{the same source}.  There is no second implementation to keep in
step.

This has a practical cost that shapes the code: branches on floating-point
comparisons do not tape correctly, so environment models are written
branch-free where possible and otherwise use conditional-expression wrappers
(\code{Environment::detail::MathWrappers}).  Table lookups are interpolated
with AD-transparent arithmetic rather than with index-selecting control flow.

\subsection{Policies checked by concepts}

A vehicle is a composition of four policies.  Their interfaces are C++20
concepts, and---this is the load-bearing detail---each concept is required to
hold for \emph{both} scalar types:

\begin{lstlisting}[caption={Policy concepts are instantiated for \code{double} and for the AD type, so a model that is not differentiable is a compile error.}, label={lst:concepts}]
template<class P, class S>
concept GravityPolicyFor =
    requires(const P& p, const Lie::SE3<S>& g, S mass) {
        { p(g, mass) } -> std::same_as<Spatial::Wrench<S>>;
    };

template<class P>
concept GravityPolicy =
       detail::GravityPolicyFor<P, double>
    && detail::GravityPolicyFor<P, CppAD::AD<double>>;
\end{lstlisting}

A policy that inadvertently uses \code{std::sqrt} on a raw \code{double}, or
takes a branch that the tape cannot record, fails the concept check at the
point of use with a readable diagnostic---rather than producing a silently
wrong Jacobian at run time.

\subsection{Predictable allocation}

The stage system, the residual, and the Newton Jacobian are fixed-size Eigen
objects and are therefore stack-resident.  No heap allocation occurs inside
the integration hot path other than the vector round-trips mandated by the
CppAD tape interface.  This is a prerequisite for the real-time and
hardware-in-the-loop use the library is intended to support.

% ---------------------------------------------------------------------------
\section{Architecture}
\label{sec:arch}
% ---------------------------------------------------------------------------

Aetherion is a header-dominant library of roughly \num{40000} lines across
the modules in Table~\ref{tab:modules}.  Dependencies are deliberately few:
Eigen for linear algebra, CppAD for differentiation, nlohmann/json for
configuration, Catch2 for tests, and fmu4cpp for FMU export.

\begin{table}[h]
  \centering
  \caption{Module map.  Namespaces mirror directory structure.}
  \label{tab:modules}
  \small
  \begin{tabular}{lll}
    \toprule
    Module & Namespace & Responsibility \\
    \midrule
    Spatial        & \code{Aetherion::Spatial}      & Twists, wrenches, spatial inertia, $\ad$/$\ad^*$, transforms \\
    ODE/RKMK/Lie   & \code{...::ODE::RKMK::Lie}     & $\SO$, $\SE$, $\Exp$, $\dexp^{-1}$, adjoints \\
    ODE/RKMK/Core  & \code{...::ODE::RKMK::Core}    & Butcher tableaux, stage residuals, Newton solve, AD Jacobian \\
    ODE/RKMK/Integrators & \code{...::Integrators}  & Radau IIA and explicit RK4 RKMK steppers \\
    RigidBody      & \code{...::RigidBody}          & State, spatial inertia assembly, vector field, 6-DOF stepper \\
    Environment    & \code{...::Environment}        & US1976 atmosphere, WGS84 gravity and constants, wind models \\
    Aerodynamics   & \code{...::Aerodynamics}       & Aerodynamic angles, Mach number, force/moment assembly \\
    Coordinate     & \code{...::Coordinate}         & ECI/ECEF/geodetic/NED conversions \\
    FlightDynamics & \code{...::FlightDynamics}     & Gravity/aero/propulsion/mass policies, trim solver \\
    Serialization  & \code{...::Serialization}      & DAVE-ML model reader, JSON configuration \\
    Simulation     & \code{...::Simulation}         & Application template method, simulator interface, CSV snapshots \\
    FMU            & ---                            & FMI 2.0/3.0 export wrappers \\
    Examples       & \code{...::Examples}           & The sixteen verification scenarios \\
    \bottomrule
  \end{tabular}
\end{table}

\subsection{Assembling a vehicle}

The composition is visible in the type.  A vehicle's right-hand side is a
\code{VectorField} parameterised by its four policies; the stepper is a
\code{SixDoFStepper} parameterised by that vector field:

\begin{lstlisting}[caption={A complete vehicle definition.  The Newton--Euler equation, the SE(3) kinematics, the implicit stage system, and the AD Jacobian are all supplied by the library.}, label={lst:vehicle}]
using RocketVF = RigidBody::VectorField<
        RocketGravityPolicy,     // w_grav(g, m)
        RocketAeroPolicy,        // w_aero(g, nu_B, m, t)
        RocketPropulsionPolicy,  // w_prop(g, nu_B, m, t)
        RocketMassPolicy>;       // mdot(t, m)

using Stepper = RigidBody::SixDoFStepper<RocketVF>;   // Radau IIA RKMK, default

Stepper stepper{ inertialParameters, newtonOptions };

auto result = stepper.step(t, state, h);
if (result.converged)
    state = Stepper::unpack(result);
\end{lstlisting}

Substituting a different integrator is a third template argument; substituting
a different vehicle is four different policy types.  Neither requires touching
the solver.

\subsection{Application scaffolding}

Executable scenarios derive from \code{Simulation::Application}, a template
method that owns the run loop, argument parsing, logging, and CSV output.  A
scenario supplies four hooks---\code{prepareSimulation}, \code{writeInitialSnapshot},
\code{stepAndRecord}, \code{logFinalSummary}---and the output column schema is
selected at compile time through \code{SnapshotTraits}, which offers a
31-column layout matching the NASA reference files exactly (so that comparison
requires no column mapping) and a 38-column extended layout that additionally
records the ECI velocity and the body-to-ECI quaternion.

% ---------------------------------------------------------------------------
\section{Numerical Core}
\label{sec:numerics}
% ---------------------------------------------------------------------------

\subsection{RKMK on $\SE \times \R^{7}$}

Runge--Kutta--Munthe-Kaas methods \citep{muntheKaas1998,iserles2000liegroup}
perform all Runge--Kutta stages inside the Lie algebra---a vector space---and
map back to the group with the exponential.  Aetherion treats the pose as a
single $\SE$ element rather than as a split $\SO \times \R^3$ product, so the
screw coupling between rotation and translation is retained.  Two integrators
are provided: a three-stage Radau~IIA scheme (order~5, stiffly accurate,
solved by Newton iteration with an exact CppAD Jacobian) and an explicit
four-stage RK4 scheme for cheap or non-stiff cases.  Both share the same stage
residual formulation and the same $\Exp$/$\dexp^{-1}$ primitives.

\begin{remark}[Trivialisation convention]
\label{rem:trivialisation}
The kinematic ODE is written in left-trivialised (body) form,
$\dot g = g \cdot \hat\xi$.  Writing $g(t) = g_0\Exp(u(t))$ then gives
$u̇ = \dexp^{-1}_{-u}(\xi)$, so the RKMK stage correction must be evaluated at
$-\eta_i$.  Evaluating it at $+\eta_i$---the right-trivialised convention,
$\dot g = \hat\xi \cdot g$, and the form in which the series is most often
tabulated---flips the sign of the $O(\ad)$ Bernoulli term and silently caps
the method at order~2 for \emph{any} Butcher tableau.  The failure is
invisible in the usual smoke tests: constant-twist motion is still integrated
exactly, and the residual still converges.  It is detected only by a
convergence study on a problem with time-varying twist.  The companion paper
\citep{tuncer2027rkmk} develops this point; it is repeated here because it is
a defect that a reader implementing RKMK is likely to reproduce.
\end{remark}

\subsection{Measured convergence order}

Order is verified rather than asserted.  The test problem is a torque-free
asymmetric rigid body ($I_{xx}, I_{yy}, I_{zz}$ all distinct) with a non-zero
initial body velocity, so that Euler's equations are nonlinear, the body twist
is genuinely time-varying, and the translational screw coupling is exercised.
Errors are measured against a reference computed at a step size 50 times
smaller than the finest tested step, cross-validated between the two
independent integrators.  Table~\ref{tab:order} reports the outcome.

\begin{table}[h]
  \centering
  \caption{Empirical convergence order at $T = \SI{2}{\second}$, torque-free
    asymmetric body.  Errors are $\norm{R_{\text{num}} - R_{\text{ref}}}_F$
    and $\norm{\mathbf{p}_{\text{num}} - \mathbf{p}_{\text{ref}}}_2$.  Mean
    log--log slope over the five step sizes.}
  \label{tab:order}
  \small
  \begin{tabular}{lcc}
    \toprule
    Integrator & Attitude order & Position order \\
    \midrule
    Radau IIA RKMK (3-stage, order 5)  & 5.00 & 5.00 \\
    Explicit RK4 RKMK (4-stage, order 4) & 3.91 & 3.98 \\
    \bottomrule
  \end{tabular}
\end{table}

A companion property is that both integrators reproduce constant-twist motion
\emph{exactly}, at every step size: a constant body twist generates a
one-parameter subgroup, and the RKMK stage system collapses onto the
exponential of a constant algebra element.  This is the defining behaviour of
a Lie-group method, and it is asserted as a regression test.

\subsection{Attitude constraint}

Over a \SI{600}{\second} tumbling trajectory the orthonormality residual
$\norm{R^\top R - I}_F$ of the SE(3) RKMK pose remains at the round-off
level and shows no secular growth, because the rotation factor is a product of
exact $\SO$ elements.  A conventional RK4 integrator applied to an
unconstrained 13-component quaternion state accumulates a residual many orders
of magnitude larger over the same trajectory.  Per-step re-normalisation of
that quaternion does restore the constraint---the distinction is therefore not
\emph{whether} orthonormality is achieved but \emph{how}: re-normalisation is a
projection applied outside the integrator, which perturbs the solution and is
not represented in any propagation Jacobian derived from the solver.
Section~\ref{sec:roadmap} returns to why this matters for filtering.

% ---------------------------------------------------------------------------
\section{Vehicle and Environment Models}
\label{sec:models}
% ---------------------------------------------------------------------------

\textbf{Atmosphere.}  The US Standard Atmosphere 1976 is implemented over the
full layered profile, returning temperature, pressure, density, and speed of
sound, in a branch-free form suitable for differentiation.

\textbf{Gravity.}  WGS84 constants and an oblateness ($J_2$) gravity model are
provided, together with the full set of ECI/ECEF/geodetic/NED conversions.
Round-trip conversion consistency is unit-tested over randomised inputs.

\textbf{Aerodynamics.}  Coefficients are read from DAVE-ML
\citep{jackson2002daveml} model files---the same interchange format used by
the NASA reference cases---and interpolated in the model's independent
variables.  Air-relative velocity accounts for the rotating atmosphere, so the
aerodynamic angles are computed against the moving air mass rather than
against the inertial frame.

\textbf{Propulsion and variable mass.}  Thrust and mass flow are supplied by a
propulsion policy and a mass policy; staging is handled by event-triggered
replacement of the active policies, with per-stage propellant tracking and a
centre-of-gravity offset that travels as propellant is consumed.  The spatial
inertia is rebuilt from the current fuel state on each step.

\textbf{Trim.}  A two-stage Newton--Raphson trim solver finds the
$(\alpha, \delta_e, \text{throttle})$ that balances a straight-and-level
condition, using an exact CppAD Jacobian for the longitudinal $2\times2$
subproblem and a bisection inversion of the propulsion map for thrust.  The
solver applies the aerodynamic-reference-centre to centre-of-gravity moment
transfer explicitly, which is necessary to reproduce the NASA F-16 trim
condition.

% ---------------------------------------------------------------------------
\section{Verification}
\label{sec:verification}
% ---------------------------------------------------------------------------

\subsection{Method}

NASA Technical Memorandum TM-2015-218675 \citep{NASA_TM_2015_218675} defines a
set of atmospheric and space flight check cases and publishes reference
trajectories generated by independently developed NASA simulation frameworks.
It is, for this purpose, close to an ideal verification target: the vehicle
models are fully specified, the reference output is machine-readable, and the
cases are graded from analytically checkable point masses up to a full
launch-vehicle ascent.

Aetherion implements sixteen of these cases as executable examples.  Each
writes a CSV file whose columns match the reference file exactly, and the two
are compared column by column with a scripted differencing tool.  The examples
are not toy drivers: they are the same code paths a user would write, and they
are built and run in continuous integration.

\subsection{Results}

Table~\ref{tab:verification} summarises agreement with the NASA reference for
each scenario.

\begin{table}[h]
  \centering
  \caption{Verification against NASA TM-2015-218675 reference trajectories.
    \TODO{Regenerate all entries after the trivialisation fix of
    Remark~\ref{rem:trivialisation}; the numbers below are carried over from
    the pre-fix documentation and must be re-measured before submission.}}
  \label{tab:verification}
  \small
  \begin{tabular}{clll}
    \toprule
    Case & Scenario & Aetherion example & Agreement with reference \\
    \midrule
     1 & Dragless sphere              & \code{DraglessSphere}            & altitude $< \SI{0.53}{\metre}$; $g$ rel.\ err.\ $< 9\times10^{-8}$ \\
     2 & Tumbling brick, undamped     & \code{TumblingBrickNoDamping}    & altitude $< \SI{0.6}{\metre}$; speed $< \SI{0.021}{\metre\per\second}$ \\
     3 & Tumbling brick, damped       & \code{TumblingBrickWithDamping}  & as case 2 translationally \\
     6 & Sphere with drag             & \code{SphereWithAtmosphericDrag} & altitude $< \SI{0.45}{\metre}$; TAS $< \SI{0.02}{\metre\per\second}$ \\
     7 & Dropped sphere, steady wind  & \code{DroppedSphereSteadyWind}   & altitude \SI{0.444}{\metre} (\SI{0.009}{\percent}); TAS $< \SI{0.001}{\percent}$ \\
     8 & Dropped sphere, wind shear   & \code{DroppedSphere2DWindShear}  & \TODO{measure} \\
     9 & Eastward cannonball          & \code{EastwardCannonball}        & \TODO{measure} \\
    10 & Northward cannonball         & \code{NorthwardCannonball}       & \TODO{measure} \\
    11 & F-16 subsonic trim           & \code{F16SteadyFlight}           & $\alpha$ within \SI{0.013}{\degree} of reference \\
    12 & F-16 supersonic trim         & \code{F16SupersonicTrim}         & trim within \SI{0.004}{\degree} \\
    13.1 & F-16 altitude change       & \code{F16AltitudeChange}         & \SI{0.4}{\foot} altitude; \SI{0.007}{\degree} pitch at \SI{20}{\second} \\
    13.2 & F-16 airspeed change       & \code{F16AirspeedChange}         & final KEAS $< \SI{0.01}{\knot}$; altitude within \SI{0.3}{\foot} \\
    13.3 & F-16 heading change        & \code{F16HeadingChange}          & final yaw within \SI{0.02}{\degree} \\
    13.4 & F-16 lateral side step     & \code{F16LateralSideStep}        & course within \SI{0.28}{\degree}; altitude within \SI{0.6}{\foot} \\
    15 & F-16 polar circumnavigation  & \code{CircleNorthPole}           & latitude $< \SI{0.0003}{\degree}$ ($<\SI{20}{\metre}$); roll $< \SI{0.1}{\degree}$ \\
    16 & F-16 equatorial circuit      & \code{CircleEquatorDateLine}     & latitude $< \SI{0.0006}{\degree}$; longitude $< \SI{0.002}{\degree}$ \\
    17 & Two-stage rocket to orbit    & \code{TwoStageRocket}            & altitude $-\SI{0.87}{\percent}$; speed $+\SI{0.14}{\percent}$ at \SI{200}{\second} \\
    \bottomrule
  \end{tabular}
\end{table}

\subsection{The launch-vehicle case}

Scenario 17 is the most demanding case and the one that exercises the whole
library at once: variable mass, travelling centre of gravity, supersonic and
hypersonic aerodynamic tables, staging, and \SI{200}{\second} of ascent from
sea level to \SI{232}{\kilo\metre}.  Through the first-stage burn the altitude
error against the NASA Sim~06 reference stays below \SI{0.15}{\percent} and the
speed error below \SI{1.5}{\metre\per\second}; the offset is unchanged through
the coast phase, confirming that the spatial inertia and gravity model are
consistent in ballistic flight; and at second-stage ignition the two
simulations are still within \SI{0.16}{\percent}.

\TODO{Insert Figure: Scenario 17 altitude and speed vs.\ NASA Sim 06, with
phase markers at S1 burnout, S2 ignition, and end of simulation.
Source: \code{doc/\_static/atmos17/}.}

Two known modelling differences account for most of the residual disagreement
at the end of the trajectory and are stated here rather than absorbed into a
tolerance.  First, the variable-mass Newton--Euler equation strictly carries an
inertia-relief term $-\dot{\IB}\nuB$ that Aetherion currently omits,
approximating it instead by rebuilding $\IB$ from the fuel state each step;
at the recommended $\Delta t = \SI{0.001}{\second}$ the per-step inertia change
is negligible against $I_{yy}$, but the approximation is not exact.  Second,
the pitch history during coast decays at a slightly different rate than the
reference, which shifts the effective second-stage ignition attitude.

% ---------------------------------------------------------------------------
\section{Extending the Library}
\label{sec:extensibility}
% ---------------------------------------------------------------------------

Adding a vehicle means writing policies, not editing the solver.  A gravity
policy is representative: it is a callable returning a body-frame wrench,
templated on the scalar type.

\begin{lstlisting}[caption={A user-supplied policy.  Writing \code{operator()} as a template over \code{S} is the entire requirement; the concept in Listing~\ref{lst:concepts} then admits it for both the simulation and the differentiated path.}, label={lst:policy}]
struct MyGravityPolicy {
    template<class S>
    Spatial::Wrench<S> operator()(const Lie::SE3<S>& g, S mass) const
    {
        const Eigen::Matrix<S,3,1> g_eci  = gravityField(g.p);   // AD-transparent
        const Eigen::Matrix<S,3,1> f_body = g.R.transpose() * (mass * g_eci);

        Spatial::Wrench<S> w{};
        w.f.template head<3>().setZero();     // moment
        w.f.template tail<3>() = f_body;      // force
        return w;
    }
};

static_assert(FlightDynamics::GravityPolicy<MyGravityPolicy>);
\end{lstlisting}

The \code{static\_assert} is the recommended idiom: it reports a violation at
the definition site, before the type is buried inside an integrator
instantiation.  Because the concept quantifies over both scalar types, a policy
that compiles for \code{double} but not for the AD type is rejected here rather
than producing a wrong Jacobian later.

The same pattern extends to the output schema (specialise
\code{SnapshotTraits} for a new column layout), to the integrator (any type
satisfying \code{IntegratorFor} can replace Radau~IIA), and to the Lie group
itself (the \code{LieGroup} concept requires only \code{Identity}, composition,
\code{Exp}, and \code{dexp\_inv}, so a different manifold can be substituted
without touching the stage machinery).

% ---------------------------------------------------------------------------
\section{Interoperability}
\label{sec:interop}
% ---------------------------------------------------------------------------

\textbf{FMI/FMU export.}  Plants are exported as FMI 2.0 and 3.0 functional
mock-up units through fmu4cpp \citep{fmu4cpp}, so an Aetherion vehicle can be
dropped into any FMI-compliant co-simulation environment.  Exported units
include the dragless sphere, the F-16 plant, an F-16 autopilot, and the
two-stage rocket.  FMU round-trip behaviour is smoke-tested in continuous
integration against the ecos co-simulation engine \citep{ecos}, and Simulink
import is exercised by a dedicated workflow.

\textbf{Model ingestion.}  Aerodynamic, inertia, propulsion, and control model
data are read from DAVE-ML files, the same format the NASA reference cases are
distributed in, so a check case can be reproduced without transcription.

\textbf{Configuration and output.}  Initial conditions and simulation
parameters are JSON; output is CSV in a column layout selectable at compile
time, including one that matches the NASA reference files exactly.

% ---------------------------------------------------------------------------
\section{Software Quality}
\label{sec:quality}
% ---------------------------------------------------------------------------

For a library whose purpose is to be trusted numerically, process is part of
the contribution.  The code base follows a practical subset of the JSF~AV C++
coding standard \citep{jsfav2005}, with deviations documented individually
rather than waived silently.  Table~\ref{tab:ci} lists the gates that run on
every pull request.

\begin{table}[h]
  \centering
  \caption{Continuous integration gates.}
  \label{tab:ci}
  \small
  \begin{tabular}{ll}
    \toprule
    Gate & Purpose \\
    \midrule
    Linux and Windows builds        & Portability across GCC/Clang and MSVC \\
    Catch2 test suite               & $\approx$\num{590} test cases in 69 translation units \\
    AddressSanitizer / UBSanitizer  & Memory and undefined-behaviour detection \\
    Clang-Tidy                      & Static analysis, coding-standard enforcement \\
    IWYU                            & Include hygiene \\
    Clang-Format / CMake-Format / CMake-Lint & Formatting and build-script consistency \\
    Metrix++                        & Cyclomatic complexity budget \\
    Codecov                         & Line coverage tracking \\
    Simulink workflow               & FMU import round trip \\
    Doxygen + Sphinx                & API reference and narrative documentation, published to GitHub Pages \\
    \bottomrule
  \end{tabular}
\end{table}

The test suite is deliberately weighted towards the numerical core: the Lie
group operations, the stage residual, the Newton solve, the spatial algebra
identities, and the coordinate conversions are each tested against closed-form
or round-trip references, independently of the end-to-end scenario
comparisons.  The convergence-order and constraint-drift studies reported in
Section~\ref{sec:numerics} are themselves regression tests, and they emit the
data files from which the figures in this paper are generated---so a change
that degrades the integrator fails the build rather than quietly degrading a
published claim.

% ---------------------------------------------------------------------------
\section{Limitations and Roadmap}
\label{sec:roadmap}
% ---------------------------------------------------------------------------

Aetherion is at version~0.11 and is explicitly pre-1.0; the following are known
gaps rather than design positions.

\begin{itemize}[leftmargin=*]
  \item \textbf{Manifold state estimation.}  The library is architected so that
    the AD Jacobian of the simulated dynamics can feed an extended Kalman
    filter on the same product manifold---this is the principal motivation for
    the constraint-preservation and AD-transparency properties above---but the
    filter itself is currently a stub.  Implementing it is the next major
    piece of work.
  \item \textbf{Inertia relief.}  The $-\dot{\IB}\nuB$ term in the
    variable-mass Newton--Euler equation is approximated by per-step rebuilding
    of $\IB$ rather than carried explicitly.
  \item \textbf{Step-size control.}  The integrators are fixed-step.  An
    embedded error estimate for Radau~IIA and adaptive stepping are planned.
  \item \textbf{Single rigid body.}  Articulated multibody dynamics are not
    implemented, although the spatial-algebra foundation is the standard one
    for it.
  \item \textbf{Aerodynamic model coverage.}  DAVE-ML ingestion covers the
    subset exercised by the NASA check cases and is not a complete
    implementation of the specification.
\end{itemize}

% ---------------------------------------------------------------------------
\section{Conclusions}
\label{sec:conclusions}
% ---------------------------------------------------------------------------

Aetherion demonstrates that the properties usually traded off against each
other in flight dynamics software---manifold correctness, differentiability,
extensibility, and verifiability---can be obtained simultaneously if they are
expressed in the type system rather than in convention.  The state is declared
on $\SE \times \R^7$ and advanced by integrators that keep it there, with
convergence orders measured at 5.00 and 3.9 rather than assumed.  Every model
is written once and evaluates identically under \code{double} and under an AD
type, with the requirement enforced by concepts that quantify over both.
Vehicles are compositions of four policies.  And the whole is verified against
sixteen published NASA reference trajectories, from analytic point masses to a
two-stage ascent to \SI{232}{\kilo\metre}, with the verification studies
wired into continuous integration.

The library is open source under the MIT licence and is offered as a
foundation for work in trajectory optimisation, manifold-consistent state
estimation, and control law verification.

% ---------------------------------------------------------------------------
\section*{Availability}
% ---------------------------------------------------------------------------

Source code, documentation, and the verification scenarios are available at
\url{https://github.com/onurtuncer/Aetherion} under the MIT licence.  API
reference and narrative documentation are published at
\url{https://onurtuncer.github.io/Aetherion/}.

% ---------------------------------------------------------------------------
\section*{Acknowledgements}
% ---------------------------------------------------------------------------

The author thanks the NASA Engineering and Safety Center (NESC) for making the
reference simulation data from TM-2015-218675 publicly available.
\TODO{Add TÜBİTAK grant number if applicable.}

% ---------------------------------------------------------------------------
% Bibliography
% ---------------------------------------------------------------------------
\bibliographystyle{unsrtnat}
\bibliography{references}

\end{document}
